\chapter{Reduction to a four-arc step function}
\label{chap:stepreduction}

% archon:covers Gluck/Euclidean/StepReduction.lean

This chapter carries out the first move of Gluck's argument (\S5--6 of the
source): replace the general continuous curvature function $\kappa$ by a
two-valued four-arc \emph{step} function $\kappa_0$ taking the values $a, b, a, b$
(with $0 < a < b$ from \cref{lem:abab_values}), and justify that solving the
closure problem for the step family solves it for $\kappa$.

\medskip

\noindent\emph{A deliberate departure from the source.} The source asserts that
two curvature functions that are $\varepsilon$-close \emph{in measure} produce
curves that are \emph{$C^1$-close}, and runs the winding argument on that basis.
That implication is \textbf{false} as stated: on the small bad set (of measure
$<\varepsilon$) where $\kappa\circ h_1$ and $\kappa_0$ disagree, the difference of
the weights $1/\kappa$ can be large, so the velocities $e^{i\theta}/\kappa$ are
\emph{not} uniformly close and the curves are not $C^1$-close. We therefore do
\emph{not} use $C^1$-closeness. What the argument actually needs --- and all it
needs --- is that the \emph{error vector} $E(\rho) = \int_0^{2\pi}
e^{i\theta}\rho(\theta)\,d\theta$ (\cref{def:error_vector}) depends continuously,
indeed Lipschitz-continuously, on the weight $\rho = 1/\kappa$ in the
$L^1$/measure sense (\cref{lem:error_vector_lipschitz}). This is the rigorous
substitute for the source's $C^1$-closeness claim, and it is what makes the
robust winding argument transfer from $\kappa_0$ to $\kappa$.

% SOURCE: [DeTurck-Gluck], §5, p. 11, lines 197-207 (read from references/deturck-gluck-fourvertex.txt)
% SOURCE QUOTE: "The first step is to replace κ by a simpler curvature function.
% Using the hypothesis that κ has at least two local maxima and two local
% minima, we find positive numbers 0 < a < b such that κ takes the values a , b ,
% a , b at four points in order around the circle. We then find a preliminary
% diffeomorphism h1 of the circle so that the function κ ◦ h1 is ε-close in
% measure to a step function κ0 with values a , b , a , b on four successive arcs
% of length π/2 each, in the sense that κ ◦ h1 is within ε of κ0 on almost all of
% S1 , except on a set of measure < ε ."
\begin{definition}

\leanok
[Four-arc step curvature]
  \label{def:four_arc_step_function}
  \lean{Gluck.stepCurvature}
  \uses{def:is_curvature_function}
  \textit{Source: DeTurck--Gluck, §5 (step function $\kappa_0$ with values $a,b,a,b$).}
  Given $0 < a < b$ and four breakpoint angles $0 \le \theta_1 < \theta_2 <
  \theta_3 < \theta_4 < 2\pi$, the \emph{four-arc step curvature} $\kappa_0 :
  \mathbb{R} \to \mathbb{R}$ is the $2\pi$-periodic function taking the value $b$
  on the arcs $[\theta_1,\theta_2)$ and $[\theta_3,\theta_4)$ and the value $a$ on
  the complementary arcs $[\theta_4, \theta_1+2\pi)$ and $[\theta_2,\theta_3)$, so
  that, read cyclically, $\kappa_0$ takes the values $a, b, a, b$ in succession.
  It is $2\pi$-periodic and strictly positive ($\ge a > 0$) but \emph{not}
  continuous: it jumps at the four breakpoints. The associated weight is
  $\rho_0 = 1/\kappa_0$, a step function with values $1/a, 1/b, 1/a, 1/b$.

  In the canonical configuration the four breakpoints are equally spaced
  ($\theta_k = (k-1)\pi/2$), giving four successive arcs of length $\pi/2$ as in
  the source; the breakpoints are left as free parameters because the winding
  argument of \cref{chap:winding} varies them over a disk of configurations.
\end{definition}

\begin{lemma}

\leanok
[Step curvature is $2\pi$-periodic]
  \label{lem:step_curvature_periodic}
  \lean{Gluck.stepCurvature_periodic}
  \uses{def:four_arc_step_function}
  The four-arc step curvature $\kappa_0$ (\cref{def:four_arc_step_function}) is
  $2\pi$-periodic: $\kappa_0(\theta + 2\pi) = \kappa_0(\theta)$ for all $\theta$.
\end{lemma}
\begin{proof}
  \uses{def:four_arc_step_function}
  \leanok
  $\kappa_0$ is defined through the reduction of $\theta$ to its representative in
  $[\theta_1, \theta_1 + 2\pi)$, which is invariant under adding $2\pi$; hence the
  value is unchanged.
\end{proof}

% Project-bespoke: elementary unfolding of the step function at the canonical
% breakpoints; no external source.
\begin{lemma}\leanok
[Canonical step-curvature values on the four arcs]
  \label{lem:step_curvature_canonical_values}
  % Lean (private, not pinned): stepCurvature_canonical_values
  \uses{def:four_arc_step_function}
  Project-bespoke. The canonical four-arc step curvature with breakpoints
  $0 < \pi/2 < \pi < 3\pi/2$ and values arranged so that it reads $a, b, a, b$
  cyclically takes its constant values on the four half-open arcs of $[0,2\pi)$:
  it equals $a$ on $[0,\pi/2)$ and on $[\pi,3\pi/2)$, and $b$ on $[\pi/2,\pi)$ and
  on $[3\pi/2,2\pi)$.
\end{lemma}
\begin{proof}
  \uses{def:four_arc_step_function}
  \leanok
  Fix $\theta$ in one of the four arcs. Since $0\le\theta<2\pi$, reducing $\theta$
  to its representative in $[0,2\pi)$ leaves it unchanged, so evaluating
  $\kappa_0$ at $\theta$ amounts to deciding which of the two indicator branches
  of \cref{def:four_arc_step_function} contains $\theta$. On $[0,\pi/2)$ and on
  $[\pi,3\pi/2)$ the selection picks the constant $a$; on $[\pi/2,\pi)$ and on
  $[3\pi/2,2\pi)$ it picks the constant $b$. Each case is a direct comparison of
  $\theta$ against the breakpoints $0,\pi/2,\pi,3\pi/2$.
\end{proof}

\begin{remark}
  $\kappa_0$ is discontinuous, so it is \emph{not} a curvature function in the
  sense of \cref{def:is_curvature_function} (which requires continuity). This is
  harmless: every statement below about $\kappa_0$ uses only that the weight
  $\rho_0 = 1/\kappa_0$ is bounded and (Lebesgue-)integrable on $[0,2\pi]$, which
  it is. The error vector $E(\rho_0)$ is defined by the same integral
  \cref{def:error_vector}, and its explicit value is computed in
  \cref{chap:bicircle}.
\end{remark}

% Project-bespoke: the L¹-Lipschitz bound on E is the rigorous replacement for
% the source's (false) "ε-close in measure ⇒ C¹-close" claim. No external source
% asserts this bound; it is elementary and stands on its proof.
\begin{lemma}

\leanok
[Error vector is $L^1$-Lipschitz in the weight]
  \label{lem:error_vector_lipschitz}
  \lean{Gluck.dist_errorVector_le}
  \uses{def:error_vector}
  For any two integrable weights $\rho, \rho' : \mathbb{R} \to \mathbb{R}$ on
  $[0,2\pi]$,
  \[
    \bigl\lVert E(\rho) - E(\rho')\bigr\rVert
      \;\le\; \int_0^{2\pi} \bigl\lvert \rho(\theta) - \rho'(\theta)\bigr\rvert \, d\theta
      \;=\; \lVert \rho - \rho'\rVert_{L^1[0,2\pi]} .
  \]
  In particular $E$ is continuous in the $L^1$ (hence the close-in-measure, for
  uniformly bounded weights) topology: if $\rho$ and $\rho'$ agree to within
  $\delta$ off a set of measure $<\varepsilon$ and are bounded by $M$, then
  $\lVert E(\rho) - E(\rho')\rVert \le 2\pi\delta + 2M\varepsilon$.
\end{lemma}
\begin{proof}
  \uses{def:error_vector}
  \leanok
  By definition $E(\rho) - E(\rho') = \int_0^{2\pi} e^{i\theta}\,(\rho(\theta) -
  \rho'(\theta))\,d\theta$ by linearity of the integral. Taking norms and using
  $\lVert\int f\rVert \le \int \lVert f\rVert$ together with $\lvert e^{i\theta}
  \rvert = 1$ gives
  \[
    \lVert E(\rho) - E(\rho')\rVert
      \le \int_0^{2\pi} \lvert e^{i\theta}\rvert \, \lvert \rho - \rho'\rvert \, d\theta
      = \int_0^{2\pi} \lvert \rho - \rho'\rvert \, d\theta .
  \]
  For the close-in-measure consequence, split $[0,2\pi]$ into the good set (where
  $\lvert\rho-\rho'\rvert \le \delta$, contributing $\le 2\pi\delta$) and the bad
  set of measure $<\varepsilon$ (where $\lvert\rho-\rho'\rvert \le 2M$,
  contributing $\le 2M\varepsilon$).
\end{proof}

\section{The preliminary reparametrization \texorpdfstring{$h_1$}{h1} (project-bespoke construction)}

\noindent\emph{Why this section exists.} The survey
(DeTurck--Gluck, §5--6) asserts the existence of a preliminary
diffeomorphism $h_1$ making $\kappa\circ h_1$ $\varepsilon$-close in measure to a
four-arc step function, but \emph{does not construct it}: it writes ``The precise
description of such a family of diffeomorphisms of the circle is given in Gluck
(1971), but we think there is enough information in the above picture for the
reader to work out the details'' (§6). We do not have Gluck (1971), so we supply a
self-contained elementary construction below. It is simpler than a literal
diffeomorphism-family argument: a single reparametrization with four ``plateaus''.

\medskip

\noindent\emph{The idea.} By \cref{lem:abab_values} the four-vertex condition
yields $0 < a < b$ and four points $c_1 < c_2 < c_3 < c_4 < c_1 + 2\pi$ in cyclic
order with $\kappa(c_1)=a$, $\kappa(c_2)=b$, $\kappa(c_3)=a$, $\kappa(c_4)=b$. A
general continuous $\kappa$ is \emph{not} close to a two-valued step function in
sup norm --- but a \emph{reparametrization} can make it close \emph{in measure}.
Build $h_1$ to be nearly constant (tiny derivative) on a large sub-arc of each
target arc $I_k = [(k-1)\pi/2,\,k\pi/2)$, mapping that sub-arc into a tiny
$\kappa$-neighbourhood $[c_k-\eta, c_k+\eta]$ of the crossing point $c_k$, and to
race through the remaining $\kappa$-domain on a short sub-arc (total $\theta$-measure
$<\delta$ per arc). Then on $\theta$-measure $\ge 2\pi - 4\delta$, $\kappa\circ h_1$
lies within the modulus of continuity $\omega(\eta)$ of the step values $a,b,a,b$,
i.e.\ $\kappa\circ h_1$ is $(\omega(\eta))$-close to $\kappa_0$ off a set of measure
$<4\delta$. Choosing $\eta,\delta$ small makes this $\varepsilon$-close in measure.

% Project-bespoke: elementary consequence of continuity; no external source.
\begin{lemma}

\leanok
[Pointwise modulus of continuity]
  \label{lem:kappa_modulus_at}
  \lean{Gluck.kappa_modulus_at}
  \uses{def:is_curvature_function}
  Project-bespoke. Let $\kappa$ be continuous and $c\in\mathbb{R}$. For every
  $\varepsilon>0$ there is $\eta>0$ such that $\lvert\kappa(t)-\kappa(c)\rvert\le
  \varepsilon$ whenever $\lvert t-c\rvert\le\eta$.
\end{lemma}
\begin{proof}

\leanok
  This is the $\varepsilon$--$\delta$ characterisation of continuity at $c$
  (\texttt{Metric.continuousAt\_iff}): pick the $\delta$ it supplies for
  $\varepsilon$ and take $\eta=\delta/2$, so $\lvert t-c\rvert\le\eta$ implies
  $\mathrm{dist}(t,c)<\delta$ and hence the bound.
\end{proof}

% Project-bespoke: the geometric core of the plateau reparametrization. No
% external source (the survey defers the explicit construction to Gluck 1971).
\begin{lemma}

[Calibrated continuous plateau density]
  \label{lem:plateau_density}
  \lean{Gluck.exists_plateau_density, Gluck.integral_affine, Gluck.tentBump, Gluck.cos_le_cos_half_shift, Gluck.tentBump_eq_zero_of_cos_le}
  % Lean (private, not pinned): tentBump_nonneg, tentBump_continuous, tentBump_periodic, cos_le_cos_half, tentBump_affine_zero, tentBump_integral_right, tentBump_integral_left, tentBump_integral_center, tentBump_integral_support, tentBump_integral_zero_of_forall, tentBump_integral_full, tentBump_integral_none, tentBump_integral_boundary, tentBump_integral_period, tentBump_integral_two_pi
  Project-bespoke. The explicit density is built from a base level plus four
  periodic unit triangular ``tent'' bumps $\mathrm{tentBump}$ (realised as
  $\arccos(\cos(\theta-\tau))$, giving $2\pi$-periodicity and continuity for free);
  the listed satellites compute their support- and cumulative integrals
  ($\mathrm{tentBump\_integral\_\ast}$) and off-support vanishing
  ($\mathrm{cos\_le\_cos\_half\ast}$, $\mathrm{tentBump\_eq\_zero\_of\_cos\_le}$),
  with $\mathrm{integral\_affine}$ for the affine pieces.
  Given crossing points $c_1<c_2<c_3<c_4<c_1+2\pi$, the offset
  $m_0=(c_1+c_4)/2-\pi$, a plateau radius $\eta>0$ small enough that
  $\eta<(c_{k+1}-c_k)/2$ for each $k$ (cyclically), and a race width $0<\delta<\pi/2$,
  there is a \emph{continuous, strictly positive, $2\pi$-periodic} speed density
  $w:\mathbb{R}\to\mathbb{R}$ with period integral $\int_0^{2\pi} w = 2\pi$ such
  that the cumulative map $h_1(\theta)=m_0+\int_0^\theta w$ satisfies, on the
  central sub-arc (``plateau'') of each target arc $I_k=[(k-1)\pi/2,\,k\pi/2)$,
  \[
    \tfrac{\delta}{2}\le\theta-(k-1)\tfrac{\pi}{2}\le\tfrac{\pi}{2}-\tfrac{\delta}{2}
    \;\Longrightarrow\; \bigl\lvert h_1(\theta)-c_k\bigr\rvert\le\eta .
  \]
\end{lemma}
\begin{proof}
  \uses{lem:kappa_modulus_at}
  Take $w$ to be the explicit continuous trapezoidal density that is large (a
  ``race'') on the two flanking sub-arcs of total length $\delta$ of each $I_k$
  and small (a ``plateau'') on the central sub-arc of length $\pi/2-\delta$,
  calibrated so the central sub-arc is carried affinely onto $[c_k-\eta,c_k+\eta]$
  and each $I_k$ is carried onto $J_k=[m_{k-1},m_k)$. Positivity holds because all
  levels are positive; $2\pi$-periodicity by matching the values at the period
  endpoints; the period integral telescopes to $\sum_k\lvert J_k\rvert=2\pi$. On
  the central plateau $h_1$ is affine with image $[c_k-\eta,c_k+\eta]$, giving the
  stated bound. The construction is elementary but lengthy (an explicit $12$-piece
  continuous piecewise-linear density with per-arc calibration); it is isolated as
  a single obligation so the surrounding \cref{lem:exists_preliminary_reparam} can
  be proved unconditionally.
\end{proof}

% Project-bespoke: elementary choice of a common small radius; no external source.
\begin{lemma}\leanok
[Existence of a common plateau radius]
  \label{lem:exists_plateau_radius}
  % Lean (private, not pinned): exists_plateau_radius
  Project-bespoke. Given four positive moduli $\eta_1,\eta_2,\eta_3,\eta_4>0$ and
  four positive gaps $g_1,g_2,g_3,g_4>0$, there is a single positive radius
  $\eta>0$ with $\eta\le\eta_i$ for every $i$ and $\eta<g_j$ for every $j$.
\end{lemma}
\begin{proof}

\leanok
  Take $\eta=\tfrac12\min(\eta_1,\eta_2,\eta_3,\eta_4,g_1,g_2,g_3,g_4)$. The
  minimum of finitely many positive numbers is positive, so $\eta>0$. As $\eta$
  is half that minimum it satisfies $\eta\le\min\le\eta_i$ for each $i$ and
  $\eta<\min\le g_j$ for each $j$ (the inequality against each $g_j$ is strict
  because halving a positive number drops it strictly below itself).
\end{proof}

% Project-bespoke: elementary Lebesgue-measure computation; no external source.
\begin{lemma}\leanok
[Total measure of the four plateau intervals]
  \label{lem:plateau_union_measure}
  % Lean (private, not pinned): plateau_union_measure
  Project-bespoke. For $0<\delta<\pi/2$ the four flanking plateau intervals
  \[
    \Bigl[\tfrac{\delta}{2},\,\tfrac{\pi}{2}-\tfrac{\delta}{2}\Bigr],\quad
    \Bigl[\tfrac{\pi}{2}+\tfrac{\delta}{2},\,\pi-\tfrac{\delta}{2}\Bigr],\quad
    \Bigl[\pi+\tfrac{\delta}{2},\,\tfrac{3\pi}{2}-\tfrac{\delta}{2}\Bigr],\quad
    \Bigl[\tfrac{3\pi}{2}+\tfrac{\delta}{2},\,2\pi-\tfrac{\delta}{2}\Bigr],
  \]
  each of length $\pi/2-\delta$ and separated by the four $\delta$-wide ``race''
  sub-arcs, have total Lebesgue measure $2\pi-4\delta$.
\end{lemma}
\begin{proof}

\leanok
  Since $0<\delta<\pi/2$, the right endpoint of each interval lies strictly below
  the left endpoint of the next (the gap between consecutive intervals is exactly
  $\delta>0$), so the four intervals are pairwise disjoint. Each has length
  $(\tfrac{\pi}{2}-\tfrac{\delta}{2})-\tfrac{\delta}{2}=\tfrac{\pi}{2}-\delta\ge0$,
  hence by additivity of Lebesgue measure on disjoint sets the union has measure
  $4\bigl(\tfrac{\pi}{2}-\delta\bigr)=2\pi-4\delta$.
\end{proof}

% Project-bespoke: not in the available survey (which defers to Gluck 1971).
% Stands on the construction below.
\begin{lemma}

\leanok
[Existence of a step-approximating reparametrization]
  \label{lem:exists_preliminary_reparam}
  \lean{Gluck.exists_preliminary_reparam}
  \uses{def:four_arc_step_function, lem:abab_values, lem:plateau_density,
        lem:kappa_modulus_at, lem:exists_plateau_radius, lem:plateau_union_measure,
        lem:step_curvature_canonical_values}
  Let $\kappa$ be a curvature function (continuous, $2\pi$-periodic, strictly
  positive) satisfying the four-vertex condition, with crossing data
  $0<a<b$, $c_1<c_2<c_3<c_4<c_1+2\pi$, $\kappa(c_1,c_2,c_3,c_4)=(a,b,a,b)$ from
  \cref{lem:abab_values}. For every $\varepsilon>0$ there is an orientation-preserving
  circle reparametrization $h_1$ (a strictly increasing continuous
  $h_1:\mathbb{R}\to\mathbb{R}$ with $h_1(\theta+2\pi)=h_1(\theta)+2\pi$) such that,
  with $\kappa_0$ the canonical four-arc step curvature
  (\cref{def:four_arc_step_function}, breakpoints $(k-1)\pi/2$, values $a,b,a,b$),
  \[
    \mathrm{vol}\bigl(\{\theta\in[0,2\pi): \lvert \kappa(h_1(\theta)) -
      \kappa_0(\theta)\rvert > \varepsilon\}\bigr) < \varepsilon .
  \]
  Moreover $h_1$ is \emph{continuously differentiable with strictly positive
  derivative}: there is a continuous $v_1:\mathbb{R}\to\mathbb{R}$ with
  $v_1(\theta)>0$ for all $\theta$ and $h_1'(\theta)=v_1(\theta)$ for all $\theta$
  (so $h_1\in C^1$ and $h_1$ is a $C^1$ diffeomorphism of the circle). This extra
  regularity is what later lets the curve be reparametrized back to realize
  $\kappa$ intrinsically (\cref{thm:gluck_converse}).
\end{lemma}
\begin{proof}
  \uses{lem:abab_values, lem:plateau_density, lem:kappa_modulus_at}
  \leanok
  Let $m_k = (c_k + c_{k+1})/2$ be the midpoints between consecutive crossings
  (cyclically, $m_4 = (c_4 + c_1 + 2\pi)/2$), so the arcs $J_k = [m_{k-1}, m_k)$
  partition the $\kappa$-circle with $c_k \in J_k$. By uniform continuity of
  $\kappa$ on the compact circle (\cref{lem:kappa_modulus_at}) pick $\eta>0$ with
  $\lvert\kappa(t)-\kappa(c_k)\rvert \le \varepsilon$ whenever
  $\lvert t - c_k\rvert \le \eta$, and small enough that
  $[c_k-\eta, c_k+\eta] \subset J_k$ for each $k$.

  \emph{The reparametrization is the antiderivative of a continuous density.}
  Apply \cref{lem:plateau_density} with these data and $m_0=(c_1+c_4)/2-\pi$ to
  obtain a \emph{continuous, strictly positive}, $2\pi$-periodic plateau density
  $w$ with $\int_0^{2\pi} w = 2\pi$ and, on the four central plateau sub-arcs of
  the target circle, the calibration
  $\bigl\lvert\, m_0 + \int_0^\theta w - c_k \,\bigr\rvert \le \eta$. Define
  \[
    h_1(\theta) \;=\; m_0 + \int_0^\theta w(s)\,ds .
  \]
  By the fundamental theorem of calculus $h_1$ is differentiable with
  $h_1'(\theta)=w(\theta)=:v_1(\theta)$; since $w$ is continuous and strictly
  positive, $h_1\in C^1$ with $v_1=w>0$ everywhere, proving the regularity claim.
  Positivity of $w$ gives $h_1$ strictly increasing; the total-integral
  normalization $\int_0^{2\pi}w=2\pi$ together with $2\pi$-periodicity of $w$ gives
  $h_1(\theta+2\pi)=h_1(\theta)+2\pi$. Thus $h_1$ is an orientation-preserving
  $C^1$ circle reparametrization.

  \emph{The measure bound.} On the union of the four plateau sub-arcs --- of total
  $\theta$-measure $\ge 2\pi-\varepsilon$ --- the calibration gives
  $h_1(\theta)=m_0+\int_0^\theta w \in [c_k-\eta, c_k+\eta]$, hence
  $\lvert\kappa(h_1(\theta)) - \kappa(c_k)\rvert \le \varepsilon$ by the choice of
  $\eta$; and on the target arc $I_k=[(k-1)\pi/2,k\pi/2)$ the step value is
  $\kappa_0=\kappa(c_k)$, so $\lvert\kappa\circ h_1 - \kappa_0\rvert \le \varepsilon$
  there. The complement (the flanking ``race'' sub-arcs where $w$ ramps between
  plateau heights) has total measure $<\varepsilon$, which bounds the exceptional
  set as claimed.
\end{proof}
